% One-round randomized prime audit bound.

\begin{theorem}[随机素数单轮审计的误接受上界]\label{thm:xi-random-prime-single-round-audit-false-accept}
\leanverified{paper\_xi\_random\_prime\_single\_round\_audit\_false\_accept}
在“综合征约束—离线审计”语义下，令 \(\{e_{\alpha}\}_{\alpha\in\mathcal I}\subset\ZZ\) 为有限族待核验的整数约束量，并假设至少存在某个 \(\alpha\) 使 \(e_\alpha\neq 0\)。
取显式有效上界 \(U\ge \max_{\alpha\in\mathcal I}|e_\alpha|\)，并令
\[
g:=\gcd_{\alpha\in\mathcal I}(e_\alpha)\ \in\ \ZZ\setminus\{0\}.
\]
从区间 \([P,2P]\) 内的素数集合中等概率抽取一素数 \(p\)。
则误接受概率满足
\[
\Pr\bigl(\forall \alpha\in\mathcal I:\ e_\alpha\equiv 0\ (\mathrm{mod}\ p)\bigr)
\ \le\
\frac{\omega(|g|)}{\pi(2P)-\pi(P)},
\qquad
\omega(|g|)\le \frac{\log|g|}{\log 2}\le \frac{N\log U}{\log 2},
\]
其中 \(N:=|\mathcal I|\)，\(\omega(m)\) 为 \(m\) 的不同素因子个数，\(\pi(\cdot)\) 为素数计数函数。
特别地，若 \(P>U\)，则误接受概率为 \(0\)。
\end{theorem}
\begin{proof}
若在模 \(p\) 下所有约束均通过，则 \(p\mid e_\alpha\) 对一切 \(\alpha\) 成立，故 \(p\mid g\)。
因此可能导致误接受的素数 \(p\in[P,2P]\) 必属于 \(g\) 的素因子集合，其个数至多 \(\omega(|g|)\)。
等概率抽取给出第一不等式。
\par
由于每个不同素因子至少贡献一个 \(2\)，有 \(2^{\omega(|g|)}\le |g|\)，从而 \(\omega(|g|)\le \log|g|/\log 2\)。
又由 \(|g|\le \prod_{\alpha}|e_\alpha|\le U^{N}\) 得 \(\log|g|\le N\log U\)，从而得到末尾界。
若 \(P>U\)，则区间内任意素数 \(p\) 不可能整除某个满足 \(0<|e_\alpha|\le U\) 的非零约束量，因此不可能误接受。证毕。
\end{proof}

\endinput
